\documentclass[12pt,a4paper]{article}

% ---------------------------------------------------------------
% Engineering report / thesis starter template
% Compile with pdfLaTeX (Overleaf: Menu -> Compiler -> pdfLaTeX)
% ---------------------------------------------------------------

\usepackage[utf8]{inputenc}
\usepackage[T1]{fontenc}
\usepackage[english]{babel}
\usepackage[margin=2.5cm]{geometry}
\usepackage{graphicx}
\usepackage{amsmath}
\usepackage{booktabs}
\usepackage{caption}
\usepackage[hidelinks]{hyperref}

\title{Title of the Report}
\author{Germ\'an Salazar}
\date{\today}

\begin{document}

\maketitle

\begin{abstract}
Short summary of the work: problem, method and main result. Keep it under 200 words.
\end{abstract}

\tableofcontents
\newpage

\section{Introduction}
State the problem and why it matters. End the section with the objective of the work.

\section{Method}
Describe the approach. Equations are numbered automatically:

\begin{equation}
  \tau = J\ddot{\theta} + b\dot{\theta} + \tau_{L}
  \label{eq:motor}
\end{equation}

Equation~\eqref{eq:motor} is referenced by its label, so numbering stays correct if
you add or remove equations.

\subsection{Materials and equipment}
List what was used. A table is often clearer than a list:

\begin{table}[h]
  \centering
  \caption{Main components.}
  \label{tab:components}
  \begin{tabular}{llr}
    \toprule
    Component & Description & Qty. \\
    \midrule
    DC motor & 12 V, geared & 1 \\
    Power stage & Two-transistor driver & 1 \\
    Controller & Microcontroller board & 1 \\
    \bottomrule
  \end{tabular}
\end{table}

Table~\ref{tab:components} is the standard way to refer to a table.

\section{Results}
Present the measurements, then interpret them. Figures go here:

% Uncomment once you add figures/result.png to the project
% \begin{figure}[h]
%   \centering
%   \includegraphics[width=0.8\textwidth]{figures/result.png}
%   \caption{Step response of the controlled plant.}
%   \label{fig:step}
% \end{figure}

\section{Conclusions}
One paragraph per finding. Do not introduce new information here.

\bibliographystyle{plain}
\bibliography{references}

\end{document}
